% GENERATED FILE. Edit canonical lessons, then run npm run generate:books.
% canonical-source-hash: fnv1a64:fa7d1b6b83c82ba3
% canonical-lessons: KA-C476-vanyajivi, KA-C476-abhayaranya, KA-C476-paksidhama, KA-C476-jivavaividhya, KA-C476-valase-hogu

\chapter{Wildlife Sanctuaries}
\label{ch:ka-wildlife-sanctuaries}
\hlchaptermodality{476}

\begin{chapteropening}
\textbf{By the end of this chapter:} \emph{I can say wildlife, a sanctuary, a bird sanctuary, biodiversity and to migrate.}

Complete the last lesson of chapter 476: I can say wildlife, a sanctuary, a bird sanctuary, biodiversity and to migrate.
\end{chapteropening}

\section[\texorpdfstring{\kn{ವನ್ಯಜೀವಿ} (vanyajīvi)}{ವನ್ಯಜೀವಿ (vanyajīvi)}]{\kn{ವನ್ಯಜೀವಿ} (vanyajīvi) — wildlife; a wild animal}

\label{lesson:KA-C476-vanyajivi}



\begin{quote}
Before the new one: say the Kannada for to fell, then the Kannada for to conserve.
\end{quote}

\subsection*{You'll want to know: \kn{ವನ್ಯಜೀವಿ}}
\textbf{\kn{ವನ್ಯಜೀವಿ}} — \emph{vanyajīvi} — \textquotedblleft{}wildlife; a wild animal\textquotedblright{}.

\begin{quote}
\kn{ವನ್ಯಜೀವಿಗಳು} \kn{ರಾತ್ರಿ} \kn{ಕಾಡಿನ} \kn{ರಸ್ತೆಯನ್ನು} \kn{ದಾಟುತ್ತವೆ}.
\end{quote}

\begin{quote}
\emph{vanyajīvigaḷu rātri kāḍina rasteyannu dāṭuttave.}
\end{quote}

\begin{quote}
Wild animals cross the forest road at night.
\end{quote}

Plural: \kn{ವನ್ಯಜೀವಿಗಳು}, \emph{vanyajīvigaḷu}.

\kn{ವನ್ಯ}, of the forest, plus \kn{ಜೀವಿ}, a living being.

\begin{cousinweb}
\textbf{\kn{ವನ್ಯಜೀವಿ}} (\emph{vanyajīvi}) is a modern compound built from Sanskrit words: \emph{vanya}, \textquotedblleft{}growing in a forest, wild\textquotedblright{} (CDIAL 11277), and \emph{jīvin}, \textquotedblleft{}living\textquotedblright{}, on \emph{jīva}, \textquotedblleft{}a living being\textquotedblright{} (CDIAL 5239).
\end{cousinweb}

\begin{grammarlens}[title={What you've built}]
One more word for getting around.
\end{grammarlens}

\subsection*{Your turn}
\begin{itemize}
\raggedright
  \item \emph{Say it:} \emph{vanyajīvi}
  \item \emph{Say it:} \emph{vanyajīvi}, once more
  \item \emph{Say it:} \emph{vanyajīvigaḷu rātri kāḍina rasteyannu dāṭuttave.} — Wild animals cross the forest road at night.
  \item \emph{Recall:} say the Kannada for to fell, then the Kannada for to conserve, then say \emph{vanyajīvi} again
\end{itemize}

\begin{tcolorbox}[breakable,colback=teal!4,colframe=teal!35!black,title={Before you move on}]
What does \kn{ವನ್ಯಜೀವಿ} mean? (\textquotedblleft{}Wildlife; a wild animal\textquotedblright{}.) Say it once more.
\end{tcolorbox}
\section[\texorpdfstring{\kn{ಅಭಯಾರಣ್ಯ} (abhayāraṇya)}{ಅಭಯಾರಣ್ಯ (abhayāraṇya)}]{\kn{ಅಭಯಾರಣ್ಯ} (abhayāraṇya) — a sanctuary; a wildlife reserve}

\label{lesson:KA-C476-abhayaranya}



\begin{quote}
Before the new one: say the Kannada for to conserve, then the Kannada for wildlife.
\end{quote}

\subsection*{You'll want to know: \kn{ಅಭಯಾರಣ್ಯ}}
\textbf{\kn{ಅಭಯಾರಣ್ಯ}} — \emph{abhayāraṇya} — \textquotedblleft{}a sanctuary; a wildlife reserve\textquotedblright{}.

\begin{quote}
\kn{ಅಭಯಾರಣ್ಯದಲ್ಲಿ} \kn{ಮರ} \kn{ಕಡಿಯುವುದು} \kn{ನಿಷೇಧ}.
\end{quote}

\begin{quote}
\emph{abhayāraṇyadalli mara kaḍiyuvudu niṣēdha.}
\end{quote}

\begin{quote}
Cutting trees in a sanctuary is banned.
\end{quote}

\kn{ಅಭಯ}, safety, plus \kn{ಅರಣ್ಯ}, forest: a safe forest.

\begin{cousinweb}
\textbf{\kn{ಅಭಯಾರಣ್ಯ}} (\emph{abhayāraṇya}) is a modern compound built from Sanskrit words: \emph{abhaya}, \textquotedblleft{}fearless; safety\textquotedblright{} (CDIAL 507), and \emph{araṇya}, \textquotedblleft{}a forest\textquotedblright{} (CDIAL 600).
\end{cousinweb}

\begin{grammarlens}[title={What you've built}]
One more word for getting around.
\end{grammarlens}

\subsection*{Your turn}
\begin{itemize}
\raggedright
  \item \emph{Say it:} \emph{abhayāraṇya}
  \item \emph{Say it:} \emph{abhayāraṇya}, once more
  \item \emph{Say it:} \emph{abhayāraṇyadalli mara kaḍiyuvudu niṣēdha.} — Cutting trees in a sanctuary is banned.
  \item \emph{Recall:} say the Kannada for to conserve, then the Kannada for wildlife, then say \emph{abhayāraṇya} again
\end{itemize}

\begin{tcolorbox}[breakable,colback=teal!4,colframe=teal!35!black,title={Before you move on}]
What does \kn{ಅಭಯಾರಣ್ಯ} mean? (\textquotedblleft{}A sanctuary; a wildlife reserve\textquotedblright{}.) Say it once more.
\end{tcolorbox}
\section[\texorpdfstring{\kn{ಪಕ್ಷಿಧಾಮ} (pakṣidhāma)}{ಪಕ್ಷಿಧಾಮ (pakṣidhāma)}]{\kn{ಪಕ್ಷಿಧಾಮ} (pakṣidhāma) — a bird sanctuary}

\label{lesson:KA-C476-paksidhama}



\begin{quote}
Before the new one: say the Kannada for wildlife, then the Kannada for a sanctuary.
\end{quote}

\subsection*{You'll want to know: \kn{ಪಕ್ಷಿಧಾಮ}}
\textbf{\kn{ಪಕ್ಷಿಧಾಮ}} — \emph{pakṣidhāma} — \textquotedblleft{}a bird sanctuary\textquotedblright{}.

\begin{quote}
\kn{ಚಳಿಗಾಲದಲ್ಲಿ} \kn{ಪಕ್ಷಿಧಾಮಕ್ಕೆ} \kn{ದೂರದ} \kn{ಹಕ್ಕಿಗಳು} \kn{ಬರುತ್ತವೆ}.
\end{quote}

\begin{quote}
\emph{caḷigāladalli pakṣidhāmakke dūrada hakkigaḷu baruttave.}
\end{quote}

\begin{quote}
In winter birds from far away come to the bird sanctuary.
\end{quote}

\kn{ಪಕ್ಷಿ}, a bird, plus \kn{ಧಾಮ}, a dwelling.

\begin{cousinweb}
\textbf{\kn{ಪಕ್ಷಿಧಾಮ}} (\emph{pakṣidhāma}) is a modern compound built from Sanskrit words: \emph{pakṣin}, \textquotedblleft{}winged; a bird\textquotedblright{} (CDIAL 7636), and \emph{dhāman}, \textquotedblleft{}a dwelling place\textquotedblright{} (Monier-Williams).
\end{cousinweb}

\begin{grammarlens}[title={What you've built}]
One more word for getting around.
\end{grammarlens}

\subsection*{Your turn}
\begin{itemize}
\raggedright
  \item \emph{Say it:} \emph{pakṣidhāma}
  \item \emph{Say it:} \emph{pakṣidhāma}, once more
  \item \emph{Say it:} \emph{caḷigāladalli pakṣidhāmakke dūrada hakkigaḷu baruttave.} — In winter birds from far away come to the bird sanctuary.
  \item \emph{Recall:} say the Kannada for wildlife, then the Kannada for a sanctuary, then say \emph{pakṣidhāma} again
\end{itemize}

\begin{tcolorbox}[breakable,colback=teal!4,colframe=teal!35!black,title={Before you move on}]
What does \kn{ಪಕ್ಷಿಧಾಮ} mean? (\textquotedblleft{}A bird sanctuary\textquotedblright{}.) Say it once more.
\end{tcolorbox}
\section[\texorpdfstring{\kn{ಜೀವವೈವಿಧ್ಯ} (jīvavaividhya)}{ಜೀವವೈವಿಧ್ಯ (jīvavaividhya)}]{\kn{ಜೀವವೈವಿಧ್ಯ} (jīvavaividhya) — biodiversity}

\label{lesson:KA-C476-jivavaividhya}



\begin{quote}
Before the new one: say the Kannada for a sanctuary, then the Kannada for a bird sanctuary.
\end{quote}

\subsection*{You'll want to know: \kn{ಜೀವವೈವಿಧ್ಯ}}
\textbf{\kn{ಜೀವವೈವಿಧ್ಯ}} — \emph{jīvavaividhya} — \textquotedblleft{}biodiversity\textquotedblright{}.

\begin{quote}
\kn{ಈ} \kn{ಕಾಡಿನ} \kn{ಜೀವವೈವಿಧ್ಯ} \kn{ತುಂಬಾ} \kn{ಶ್ರೀಮಂತ}.
\end{quote}

\begin{quote}
\emph{ī kāḍina jīvavaividhya tumbā śrīmanta.}
\end{quote}

\begin{quote}
The biodiversity of this forest is very rich.
\end{quote}

\kn{ಜೀವ}, life, plus \kn{ವೈವಿಧ್ಯ}, variety.

\begin{cousinweb}
\textbf{\kn{ಜೀವವೈವಿಧ್ಯ}} (\emph{jīvavaividhya}) is a modern compound built from Sanskrit \emph{jīva}, \textquotedblleft{}a living being, life\textquotedblright{} (CDIAL 5239), and \emph{vaividhya}, \textquotedblleft{}variety\textquotedblright{}, from \emph{vividha}, \textquotedblleft{}of many kinds\textquotedblright{}.
\end{cousinweb}

\begin{grammarlens}[title={What you've built}]
One more word for getting around.
\end{grammarlens}

\subsection*{Your turn}
\begin{itemize}
\raggedright
  \item \emph{Say it:} \emph{jīvavaividhya}
  \item \emph{Say it:} \emph{jīvavaividhya}, once more
  \item \emph{Say it:} \emph{ī kāḍina jīvavaividhya tumbā śrīmanta.} — The biodiversity of this forest is very rich.
  \item \emph{Recall:} say the Kannada for a sanctuary, then the Kannada for a bird sanctuary, then say \emph{jīvavaividhya} again
\end{itemize}

\begin{tcolorbox}[breakable,colback=teal!4,colframe=teal!35!black,title={Before you move on}]
What does \kn{ಜೀವವೈವಿಧ್ಯ} mean? (\textquotedblleft{}Biodiversity\textquotedblright{}.) Say it once more.
\end{tcolorbox}
\section[\texorpdfstring{\kn{ವಲಸೆ} \kn{ಹೋಗು} (valase hōgu)}{ವಲಸೆ ಹೋಗು (valase hōgu)}]{\kn{ವಲಸೆ} \kn{ಹೋಗು} (valase hōgu) — to migrate}

\label{lesson:KA-C476-valase-hogu}



\begin{quote}
Before the new one: say the Kannada for a bird sanctuary, then the Kannada for biodiversity.
\end{quote}

\subsection*{You'll want to know: \kn{ವಲಸೆ} \kn{ಹೋಗು}}
\textbf{\kn{ವಲಸೆ} \kn{ಹೋಗು}} — \emph{valase hōgu} — \textquotedblleft{}to migrate\textquotedblright{}.

\begin{quote}
\kn{ಕೆಲವು} \kn{ಹಕ್ಕಿಗಳು} \kn{ಚಳಿಗಾಲದಲ್ಲಿ} \kn{ದಕ್ಷಿಣಕ್ಕೆ} \kn{ವಲಸೆ} \kn{ಹೋಗುತ್ತವೆ}.
\end{quote}

\begin{quote}
\emph{kelavu hakkigaḷu caḷigāladalli dakṣiṇakke valase hōguttave.}
\end{quote}

\begin{quote}
Some birds migrate south in winter.
\end{quote}

Past (I …): \kn{ವಲಸೆ} \kn{ಹೋದೆ}, \emph{valase hōde}.

Of birds, and of people who leave home for work: \kn{ವಲಸೆ} \kn{ಕಾರ್ಮಿಕ}, a migrant worker.

\begin{cousinweb}
\textbf{\kn{ವಲಸೆ} \kn{ಹೋಗು}} (\emph{valase hōgu}) is native Dravidian: DEDR 5278 lists Kannada \emph{valase}, \textquotedblleft{}flight, removal from home, emigration\textquotedblright{}, here with \kn{ಹೋಗು}, to go.
\end{cousinweb}

\begin{grammarlens}[title={What you've built}]
One more word for getting around.
\end{grammarlens}

\subsection*{Your turn}
\begin{itemize}
\raggedright
  \item \emph{Say it:} \emph{valase hōgu}
  \item \emph{Say it:} \emph{valase hōgu}, once more
  \item \emph{Say it:} \emph{kelavu hakkigaḷu caḷigāladalli dakṣiṇakke valase hōguttave.} — Some birds migrate south in winter.
  \item \emph{Recall:} say the Kannada for a bird sanctuary, then the Kannada for biodiversity, then say \emph{valase hōgu} again
\end{itemize}

\begin{tcolorbox}[breakable,colback=teal!4,colframe=teal!35!black,title={Before you move on}]
What does \kn{ವಲಸೆ} \kn{ಹೋಗು} mean? (\textquotedblleft{}To migrate\textquotedblright{}.) Say it once more.
\end{tcolorbox}
